\documentclass{article}

\usepackage[utf8]{inputenc}     % pdfLaTeX
\usepackage[T1]{fontenc}
\usepackage[magyar]{babel}

\usepackage{amsmath}
\usepackage{amsthm}
\usepackage{amssymb}

\title{Tiszta algebra}
\author{Rovenszky Robin\thanks{Rovenszky Robin fordításával}}
\date{Utoljára frissítve: 2026. március 16.}

\newtheorem{theorem}{Tétel}
\newtheorem*{remark}{Megjegyzés}

\begin{document}

\maketitle

\section{Mi az a $\sqrt{2}$?}
Tudjuk, hogy léteznie kell, hiszen egy derékszögű háromszögben, ahol a befogók hossza $a = 1$, $b = 1$, és a közöttük lévő szög $\alpha = 90^\circ$, az átfogó hossza $\sqrt{2}$. De mi is valójában $\sqrt{2}$?\\

\begin{theorem}
    Legyen $x$ olyan szám, amelyre $x^2 = 2$. Feltehetjük, hogy $x$ kifejezhető $\frac{p}{q}$ alakban, ahol $p, q \in \mathbb{Z}$ és $\gcd(p,q) = 1$.
\end{theorem}

\begin{proof}
    Tegyük fel, hogy az állítás igaz. Ekkor
    \begin{equation}
        \frac{p^2}{q^2} = 2 \implies p^2 = 2q^2 \implies 2 \mid p^2 \implies 2 \mid p \implies 4 \mid p^2 \implies
    \end{equation}
    \begin{equation}
        \implies 4 \mid 2q^2 \implies 2 \mid q^2 \implies 2 \mid q \implies \gcd(p, q) \ge 2 \Rightarrow\!\Leftarrow
    \end{equation}
    Vagy alternatív módon befejezhettük volna a bizonyítást $p^2 = 2q^2$ segítségével. Ekkor észrevesszük, hogy $p^2$ prímtényezős felbontásában a kettes kitevője páros, míg $2q^2$ prímtényezős felbontásában a kettes kitevője páratlan (mivel $\gcd(p,q) = 1$).
\end{proof}

Ezt az eredményt általánosíthatjuk úgy, hogy minden $a \in \mathbb{Z}_{\ge 0}$ esetén, ha $x^2 = a$, és $a$ tartalmaz prímtényezőt páratlan kitevővel (azaz $a$ nem tökéletes négyzet), akkor és csak akkor $x \in \mathbb{R} \setminus \mathbb{Q}$.

\section{$\sqrt{2}$ meghatározása}
Mivel irracionális számot nem találhatunk pontosan, közelítjük. Kezdjük a következő intervallummal:
\[
1 < x < 2 \Leftrightarrow 1 < x^2 < 4
\]

\begin{theorem}
    \[
    0 < x < y \implies 0 < x^2 < y^2
    \]
\end{theorem}

\hrule
\vspace{1em}

Mielőtt bizonyítanánk a tételt, írjuk le a **rendezési axiómákat**:

\begin{enumerate}
    \item \textbf{Trikotómia axióma:} Pontosan egy állítás lehet igaz a következők közül: $a < b$, $b < a$, $a = b$.
    \item \textbf{Összeadás megőrzi a rendezést:} Ha $a < b$ és $c$ tetszőleges, akkor $a+c < b+c$.
    \item \textbf{Tranzitivitás axiómája:} Ha $a < b$ és $b < c$, akkor $a < c$.
    \item \textbf{Pozitív szorzás megőrzi a rendezést:} Ha $0 < a$ és $b < c$, akkor $ab < ac$.
    \item \textbf{1 pozitivitása:} $0 < 1$.
\end{enumerate}

Az összeadás disztributivitását (vagy hogy az összeadás megőrzi a rendezést) használva bizonyíthatjuk:
\[
x < 0 \implies 0 < -x
\]

\begin{proof}
\[
x + (-x) < 0 + (-x)
\]
\[
0 < -x
\]
\end{proof}

A szorzás disztributivitását (vagy hogy a pozitív szorzás megőrzi a rendezést) használva bizonyíthatjuk:
\[
(-1)\cdot x = -x
\]

\begin{proof}
\[
(1 + (-1)) \cdot x = 0
\]
\[
1 \cdot x + (-1) \cdot x = 0 \implies (-1)\cdot x = -x
\]
\end{proof}

Ehhez szükségünk van még a következő állításra is:
\[
0 \cdot x = 0
\]

\begin{proof}
\[
(0+0)\cdot x = 0\cdot x
\]
\[
0\cdot x + 0\cdot x = 0\cdot x
\]
\[
0\cdot x = 0
\]
\end{proof}

\hrule
\vspace{1em}

Most bizonyítsuk a 3. tételt a szorzás disztributivitásának segítségével:

\begin{proof}
\begin{equation}
0 < x < y \implies 0 < x^2 < y^2
\end{equation}

Szorzunk $x$-szel:
\[
x^2 < xy
\]

Majd (27)-et $y$-val:
\[
xy < y^2
\]

Egyesítve:
\[
x^2 < xy < y^2 \implies x^2 < y^2
\]
\end{proof}

\begin{remark}
Szintén bizonyítható, hogy $y < x < 0 \implies 0 < -x < -y$ a szorzás disztributivitása segítségével.
\end{remark}

\begin{proof}
\[
(-x) \cdot (-x) < (-y) \cdot (-x)
\]
\[
x^2 < xy < y^2
\]
\end{proof}

\end{document}